\documentclass[11pt]{article}
\usepackage[margin=1in]{geometry}
\usepackage{amsmath,amssymb,amsthm}
\usepackage{mathtools}
\usepackage{bm}
\usepackage{enumitem}
\usepackage{hyperref}

\newtheorem*{exercise}{Exercise}

\newcommand{\ket}[1]{\lvert #1 \rangle}
\newcommand{\bra}[1]{\langle #1 \rvert}
\newcommand{\braket}[2]{\langle #1 \vert #2 \rangle}
\newcommand{\F}{\mathbb{F}_2}
\newcommand{\Pn}{\mathcal{P}_n}
\newcommand{\Sop}{\mathcal{S}}

\title{ATPL 2026 --- Exercises: Stabilizers and the Clifford Group}
\author{}
\date{}

\begin{document}
\maketitle
\vspace{-2em}

\noindent
These exercises accompany Lecture 6 (Stabilizers, Clifford Algebra, and separating
Classical Computations from Quantum). To supplement the exercises, QPP contains a
code file \texttt{StabilizersProblem.hs}, with exercises for Haskell and QPP
library training, similar to the topics found here.

\medskip
\noindent
\textbf{Background reading:} Nielsen \& Chuang, Section 10.5 (Stabilizer
codes); Aaronson \& Gottesman (2004), \emph{Improved simulation of stabilizer
circuits}; Garc\'ia, Markov \& Cross (2014), \emph{On the geometry of
stabilizer states}.

\bigskip
\section*{Part 1: Stabilizers and the Pauli group}

\begin{exercise}[Closure of the Pauli group]
Recall $\mathcal{P}_1 = \{ i^l \sigma \mid \sigma \in \{I,X,Y,Z\},\ l \in \{0,1,2,3\}\}$.
\begin{enumerate}[label=(\alph*)]
    \item Using the multiplication table from the lecture (or by direct
    matrix multiplication), verify that $XY = iZ$ and $YX = -iZ$, and hence
    that $\mathcal{P}_1$ is \emph{non-abelian}.
    \item Show that every element of $\mathcal{P}_1$ has a well-defined
    inverse within $\mathcal{P}_1$, and identify the inverse of $iY$.
    \item What is $|\mathcal{P}_1|$? What is $|\Pn|$ in general? (Count the
    choices of $\sigma_k \in \{I,X,Y,Z\}$ for each of the $n$ tensor factors,
    together with the single overall phase $i^l$.)
\end{enumerate}
\end{exercise}

\begin{exercise}[Properties of the stabilizer group, done carefully]
Let $\ket{\psi}$ be a normalized $n$-qubit state and
$\Sop(\psi) = \{P \in \Pn \mid P\ket\psi = \ket\psi\}$.
\begin{enumerate}[label=(\alph*)]
    \item Show directly (do not just cite the lecture) that every
    $P \in \Sop(\psi)$ satisfies $P^2 = I$. \emph{Hint:} $P^2$ is itself an
    element of $\Pn$, hence of the form $i^l I$ for some $l$; use
    $P\ket\psi=\ket\psi$ twice.
    \item Show that $\Sop(\psi)$ is abelian. \emph{Hint:} for $P,Q \in \Pn$,
    either $PQ = QP$ or $PQ = -QP$ (this is a general fact about Pauli
    strings --- you may assume it, or prove it using the single-qubit
    multiplication table applied factor-by-factor). Rule out the
    anticommuting case using $P,Q \in \Sop(\psi)$.
    \item Why does part (a) already rule out $-I \in \Sop(\psi)$ for any
    valid state $\ket\psi$? What would it mean physically if $-I$ stabilized
    a state?
\end{enumerate}
\end{exercise}

\begin{exercise}[Computing stabilizers from states]
For each state below, find a full set of independent stabilizer generators
(i.e.\ enough elements of $\Pn$ to pin down the state uniquely, per
Theorem~1 of the lecture).
\begin{enumerate}[label=(\alph*)]
    \item $\ket{00}$
    \item $\ket{+}\ket{0}$, where $\ket{+} = \tfrac{1}{\sqrt2}(\ket0+\ket1)$
    \item The Bell state $\ket{\Phi^+} = \tfrac{1}{\sqrt2}(\ket{00}+\ket{11})$
    \item The 3-qubit GHZ state $\tfrac{1}{\sqrt2}(\ket{000}+\ket{111})$
\end{enumerate}
For (c) and (d), check explicitly that your proposed generators commute
with one another (a necessary condition, by Exercise 2b, for them to
simultaneously stabilize a state).
\end{exercise}

\begin{exercise}[From generators to state]
Consider the two-qubit operators $X\otimes X$ and $Z \otimes Z$.
\begin{enumerate}[label=(\alph*)]
    \item Verify that $X\otimes X$ and $Z\otimes Z$ commute, so that they
    could simultaneously stabilize a state.
    \item Following the method of the lecture (Heisenberg representation of
    states, slides 12--14), solve the two linear systems
    $\ket\psi = (X\otimes X)\ket\psi$ and $\ket\psi = (Z\otimes Z)\ket\psi$
    simultaneously to find the unique normalized state stabilized by both.
    Confirm it matches Exercise 3(c).
\end{enumerate}
\end{exercise}

\begin{exercise}[Counting stabilizer states, $n=1$]
By enumerating the $\pm X, \pm Y, \pm Z$ eigenstates found on lecture
slide~11, list all six 1-qubit stabilizer states explicitly (as vectors in
the computational basis). Convince yourself that no other single-qubit
state is a stabilizer state (i.e.\ has $|\Sop(\psi)|=2$, per Theorem~1 with
$n=1$).
\end{exercise}

\bigskip
\section*{Part 2: Check matrices (stabilizer tableaux)}

Recall from the lecture that a Pauli group element can be written using two
bits per qubit, $\sigma = i^{x z} X^x Z^z$, giving the encoding
$[I]=00,\ [Z]=01,\ [X]=10,\ [Y]=11$, and that a general $n$-qubit element is
represented by a $(2n+2)$-bit string $(s,c,\mathbf{x},\mathbf{z})$ via
\[
    [P] = \big[(-1)^s i^c X^{x_0}Z^{z_0}\otimes \cdots \otimes X^{x_{n-1}}Z^{z_{n-1}}\big]
        = (s,c,\mathbf{x},\mathbf{z}).
\]

\begin{exercise}[Encoding and decoding]
\begin{enumerate}[label=(\alph*)]
    \item Encode $P = -Y \otimes I \otimes Z$ as a bit string
    $(s,c,\mathbf{x},\mathbf{z})$.
    \item Decode the bit string $(s,c,\mathbf x,\mathbf z) = (0,1,101,110)$
    (for $n=3$) back into a signed tensor product of Pauli matrices.
\end{enumerate}
\end{exercise}

\begin{exercise}[Multiplication via XOR]
Let $P_1 = Y \otimes Z \otimes I$ and $P_2 = X \otimes Z \otimes X$.
\begin{enumerate}[label=(\alph*)]
    \item Encode $P_1$ and $P_2$ as bit strings, and use the rule
    \[
        [P_1 P_2] = [P_1] \oplus [P_2] \oplus \text{(sign correction)}, \qquad
        \text{sign correction} = \mathbf{z}_1 \cdot \mathbf{x}_2 \ (\mathrm{mod}\ 2)
        \text{ on the sign bit},
    \]
    to compute $[P_1 P_2]$.
    \item Decode the result back into a signed tensor product, and check it
    against direct matrix multiplication of $P_1 P_2$ qubit-by-qubit.
\end{enumerate}
\end{exercise}

\begin{exercise}[The symplectic product and commutativity]
For two single-qubit-per-factor Pauli strings with bit vectors
$(\mathbf x_1,\mathbf z_1)$ and $(\mathbf x_2,\mathbf z_2)$ (ignoring phase
bits), define the \textbf{symplectic product}
\[
    \langle (\mathbf x_1,\mathbf z_1), (\mathbf x_2,\mathbf z_2)\rangle
    \;=\; \mathbf x_1 \cdot \mathbf z_2 \;+\; \mathbf z_1 \cdot \mathbf x_2
    \pmod 2,
\]
where $\cdot$ is the usual dot product over $\F$.
\begin{enumerate}[label=(\alph*)]
    \item Show, one qubit at a time, that two single-qubit Pauli operators
    $P = i^{x_1z_1}X^{x_1}Z^{z_1}$ and $Q = i^{x_2z_2}X^{x_2}Z^{z_2}$
    commute if and only if $x_1z_2 + z_1x_2 \equiv 0 \pmod 2$.
    \emph{Hint:} check the four cases $(x_1,z_1) = 00,10,01,11$ against
    each of $(x_2,z_2)$, or use $XZ=-ZX$ together with the fact that $I$
    and any single operator commute with everything.
    \item Extend this to $n$-qubit strings: show that $P_1$ and $P_2$
    commute if and only if their symplectic product $\langle
    (\mathbf x_1,\mathbf z_1),(\mathbf x_2,\mathbf z_2)\rangle$, summed over
    all $n$ qubits, vanishes mod 2. (This is exactly the check-matrix
    version of the commutation criterion used implicitly on lecture
    slide~16.)
    \item Use this criterion to verify your commutation check from
    Exercise 4(a), and to check whether $X\otimes Z \otimes I$ and $Z
    \otimes I \otimes X$ commute.
\end{enumerate}
\end{exercise}

\begin{exercise}[Independence of generators]
Recall Theorem~1: an $n$-qubit state is a genuine stabilizer state
\emph{iff} it has $n$ \emph{independent} commuting generators. Independence
here means linear independence of the $(\mathbf x,\mathbf z)$ rows over
$\F$ (the sign/phase bits $s,c$ do not affect independence).

\noindent
For $n=3$, consider the following three candidate generators, given as
rows $(\mathbf x \mid \mathbf z)$:
\[
    P_1 = (110\mid 001), \qquad
    P_2 = (011\mid 100), \qquad
    P_3 = (101\mid 101).
\]
\begin{enumerate}[label=(\alph*)]
    \item Check pairwise commutativity using the symplectic product from
    Exercise 8.
    \item Stack the three $(\mathbf x\mid\mathbf z)$ rows into a $3\times 6$
    matrix over $\F$ and compute its rank via Gaussian elimination mod~2.
    Is $\{P_1,P_2,P_3\}$ an independent generating set?
    \item Now replace $P_3$ with $P_3' = P_1 \oplus P_2$ (bitwise XOR, i.e.\
    the tableau product of $P_1$ and $P_2$ up to sign). Explain --- without
    redoing the rank computation --- why $\{P_1,P_2,P_3'\}$ can \emph{not}
    be an independent generating set.
\end{enumerate}
\end{exercise}

\bigskip
\section*{Part 3: Clifford vs.\ non-Clifford, and minimal form}

\begin{exercise}[Verifying the conjugation table]
The lecture states that $H$, $S$ (the phase gate $P$), and $\mathrm{CNOT}$
generate the Clifford group, with (among others) $HXH^\dagger = Z$,
$HZH^\dagger = X$, $HYH^\dagger = -Y$, and $SXS^\dagger = Y$,
$SYS^\dagger = -X$, $SZS^\dagger = Z$.
\begin{enumerate}[label=(\alph*)]
    \item Verify $HXH^\dagger = Z$ and $HYH^\dagger = -Y$ by direct $2\times
    2$ matrix multiplication, using $H = \tfrac{1}{\sqrt2}
    \begin{psmallmatrix}1&1\\1&-1\end{psmallmatrix}$.
    \item Verify $SXS^\dagger = Y$ by direct matrix multiplication, using
    $S = \begin{psmallmatrix}1&0\\0&i\end{psmallmatrix}$.
    \item Both $H$ and $S$ map $\mathcal{P}_1 \to \mathcal{P}_1$ under
    conjugation (up to sign). Explain briefly why this is exactly the
    defining property of the Clifford group stated in the lecture, and why
    a generic single-qubit unitary $U(\theta) = e^{i\theta X}$ (for
    irrational $\theta/\pi$) does \emph{not} have this property.
\end{enumerate}
\end{exercise}

\begin{exercise}[Bit-level update rules]
Because $H$ and $S$ permute $\Pn$ under conjugation, they induce simple
\emph{bit-level} update rules on a stabilizer tableau row $(x,z)$ (per
qubit they act on), rather than requiring full matrix multiplication.
\begin{enumerate}[label=(\alph*)]
    \item Using the identities $HXH^\dagger=Z$, $HZH^\dagger=X$, derive the
    update rule for $H$ acting on a single-qubit tableau entry: what
    happens to $(x,z)$, and does the sign bit $s$ ever flip? (Check the
    $Y=iXZ$ case, i.e.\ $(x,z)=(1,1)$, carefully --- this is where a sign
    correction can appear.)
    \item Do the same for $S$: derive the update rule for $(x,z) \mapsto
    (x',z')$ under $S$-conjugation, again checking whether $(1,1)$
    (i.e.\ $Y$) picks up a sign.
    \item Explain why these per-qubit bit-flip/XOR rules, applied to every
    row of the tableau in parallel, reproduce the $\mathcal{O}(n)$-per-gate
    update claimed on lecture slide~28.
\end{enumerate}
\end{exercise}

\begin{exercise}[Pushing a non-Clifford gate through the Clifford frame]
Let $T = \begin{psmallmatrix}1&0\\0&e^{i\pi/4}\end{psmallmatrix}$ (the
non-Clifford $T$ gate).
\begin{enumerate}[label=(\alph*)]
    \item Compute $STS^\dagger$ explicitly. Show that it equals $T$ up to
    an overall phase (i.e.\ $STS^\dagger = e^{i\phi}T$ for some $\phi$) ---
    in this special case, commuting $T$ past the Clifford gate $S$ leaves
    it essentially unchanged.
    \item Compute $HTH^\dagger$ explicitly, and show that it is \emph{not}
    a phase multiple of $T$, but is instead a genuinely different
    non-Clifford single-qubit rotation. (You do not need to name it; just
    exhibit the matrix.)
    \item Given the one-qubit circuit $U = H\, T\, H\, T\, H$ (read
    right-to-left as applied to a ket), identify which gates are Clifford
    and which are not. Using the identity $C N C^\dagger$ (apply $C$, then
    the transformed non-Clifford gate, i.e.\ rewrite $CN = (CNC^\dagger)C$)
    repeatedly, rewrite $U$ in the form $N_1' N_2' C$ where $C$ is a single
    net Clifford operator and $N_1', N_2'$ are (possibly modified)
    non-Clifford gates coming from the two $T$'s. You do not need to
    simplify $N_1', N_2'$ beyond expressing them as conjugated $T$'s; the
    point of the exercise is the bookkeeping, i.e.\ tracking the ``Clifford
    frame'' as in lecture slide~31.
\end{enumerate}
\end{exercise}

\begin{exercise}[Why $T$-count, not gate count, dominates the cost --- optional discussion]
The Aaronson--Gottesman extension (lecture slide~4) simulates a circuit
with $M$ Clifford gates and $T$ non-Clifford gates in time
$\mathcal{O}(n^2 + nM + nT\,2^{\gamma T})$.
\begin{enumerate}[label=(\alph*)]
    \item Explain, using your bookkeeping from Exercise 12(c), why every
    Clifford gate in the circuit can be absorbed into the classically
    tracked stabilizer tableau at cost $\mathcal{O}(n)$ regardless of how
    many non-Clifford gates come later, while each non-Clifford gate
    contributes an exponential factor.
    \item If you were designing a quantum algorithm and had a choice
    between two circuits that compute the same function, one using $50$
    Clifford gates and $5$ $T$-gates, and another using $10$ Clifford gates
    and $8$ $T$-gates, which would you expect to be cheaper to
    \emph{classically pre-simulate/verify}, and why? (This is exactly the
    motivation for ``$T$-count minimization'' as a circuit-optimization
    goal in fault-tolerant quantum computing.)
\end{enumerate}
\end{exercise}

\end{document}
